 
% kuvien shiftaukset/marginaalit kalvoihin
\xdef\figYshift{-0.75cm}
\xdef\figXshift{-0.25cm}


%% itemize-ympäristöjen marginaalit, rivivälit jne.
\xdef\iilm{0.5cm}
\xdef\ilm{-0.0cm}
\xdef\isep{2mm}
%\xdef\semismall{\fontsize{11}{12}\selectfont}

%% Nuoli-implikaatio itemize-ympäristönä
\renewcommand{\implies}[2][\fontsize{11}{12}\selectfont]{\vspace{1mm}% new addition, done in 
\begin{itemize}[leftmargin=0.7cm,
			rightmargin=-0.5cm,
			itemsep=3mm] 
			#1 %28.10.2021 addition
        \item[=>] #2
    \end{itemize}
}
\newcommand{\Implies}[2][\fontsize{11}{12}\selectfont]{\vspace{1mm}% new addition, done in 
\begin{itemize}[leftmargin=0.7cm,
			rightmargin=-0.5cm,
			itemsep=3mm] 
			#1 %28.10.2021 addition
        \item[=>]<1-> #2
    \end{itemize}
}

\newcommand{\RA}{\Rightarrow} % reduces problems with ctrl+f
\newcommand{\ra}{\rightarrow} % reduces problems with ctrl+f

%\setlist{nosep,topsep=-\parskip}
\setlist[itemize]{leftmargin=-0.35cm,
			rightmargin=-0.5cm,
			itemsep=2.5mm,
			,label=\i} %3mm 
			
%https://tex.stackexchange.com/questions/533206/nested-enumerated-list-with-different-left-margins-in-latex
% Create new environments for listings
%\newlist{subitemize}{enumerate}{3}
%\setlist[subitemize]{
%  wide,
%  align=left,
%  topsep = 0in,
%  parsep = 0.24in,
%  partopsep = 0in,
%  itemsep = 0in,
%  labelsep = 0in,
%  itemindent=0.5in,
%  labelindent = -0.5in,
%  labelwidth=0.5in,
%}
%\setlist[subitemize,1]{leftmargin=0in,label=\arabic*.}
%\setlist[subitemize,2]{leftmargin=0.5in,label=(\alph*)}
%\setlist[subitemize,3]{leftmargin=1in,label=(\roman*)}

%% modify the nested itemize loops, in order to keep things simple:
%\setlist[itemize,1]{leftmargin=-0.35cm,rightmargin=-0.5cm,label=\i,itemsep=2.5mm}
\setlist[itemize,2]{leftmargin=0.5cm,rightmargin=0.4cm,itemsep=0.5mm,label=\ii}
\setlist[itemize,3]{leftmargin=1cm,itemsep=1mm,label=\iii}

%% Slide symbols:
\xdef\i{-} %*. {\symbol{"25CA}}$ ~diamond
\xdef\ii{\symbol{"25CA}} %*. {\symbol{"25CA}}$ ~diamond
\xdef\iii{\symbol{"27A4}} %*. {\symbol{"25CA}}$ ~diamond

%% SoM Slide symbols:  > , circle, -
\xdef\i{\symbol{"27A4}} %*. {\symbol{"25CA}}$ ~diamond
\xdef\ii{\symbol{"25CF}} %*. {\symbol{"25CA}}$ ~diamond
\xdef\iii{-} %*. {\symbol{"25CA}}$ ~diamond

% \global pause to slide items
%\beamerdefaultoverlayspecification{<+->}
%
%\newcommand{\appear}[1]{%
%\uncover<+->{%
%   #1
%}%
%}

%\newcommand{\appear}[1]{%
%   #1
%}

% Set up the keys.  Only the ones directly under /myparbox
% can be accepted as options to the \myparbox macro.
\pgfkeys{
 /myfootnote/.is family, /myfootnote,
 % Here are the options that a user can pass
 default/.style = {appear= .-, position = right, ref = H},
 appear/.estore in = \myappearCMD,
 Appear/.estore in = \myappearCMD,
 position/.style = {positions/#1/.get = \myfootnotePosition},
 Position/.style = {positions/#1/.get = \myfootnotePosition},
 pos/.style = {positions/#1/.get = \myfootnotePosition},
% Ref/.estore in = \myfootnoteRef,
 ref/.style = {positions/#1/.get = \myfootnoteRef},
 Ref/.style = {positions/#1/.get = \myfootnoteRef},
 % /handlers/first char syntax=true,
%  /handlers/first char syntax/the character </.initial=\mypointedmacro,
 % Here is the dictionary for positions.
 positions/.cd,
  left/.initial = south west,
  l/.initial = south west,
  L/.initial = south west,
  west/.initial = south west,
  right/.initial = south east,
  r/.initial = south east,
  R/.initial = south east,
  east/.initial = south east,
% dictionary for the reference texts:
  Harris/.initial = {Figure from R. Harris, \textit{Modern Physics} (2$^{nd}$ edition).},
  H/.initial = {Figure from R. Harris, \textit{Modern Physics} (2$^{nd}$ edition).},
  h/.initial = {R. Harris, \textit{Modern Physics} (2$^{nd}$ edition).},
  HarrisShort/.initial = {R. Harris, \textit{Modern Physics} (2$^{nd}$ edition).},
  Tipler/.initial = {Figure from P. Tipler \& R. Llewellyn, \textit{Modern Physics} (5$^{th}$ edition).},
  T/.initial = {Figure from P. Tipler \& R. Llewellyn, \textit{Modern Physics} (5$^{th}$ edition).},
  t/.initial = {P. Tipler \& R. Llewellyn, \textit{Modern Physics} (5$^{th}$ edition).},
  TiplerShort/.initial = {P. Tipler \& R. Llewellyn, \textit{Modern Physics} (5$^{th}$ edition).},
  Harris and Tipler/.initial = {Figure from R. Harris, \textit{Modern Physics} (2$^{nd}$ edition).},
  h and t/.initial = {R. Harris, \textit{Modern Physics} (2$^{nd}$ edition) and P. Tipler \& R. Llewellyn, \textit{Modern Physics} (5$^{th}$ edition).},
%  HarrisShort/.initial = R. Harris, \textit{Modern Physics} (2$^{nd}$ edition).,
%  TiplerShort/.initial = P. Tipler \& R. Llewellyn, \textit{Modern Physics} (5$^{th}$ edition).,
%  
}
%\def\mypointedmacro#1{{\pgfkeysalso{appearCMD={#1}}}}
\newcommand\CR[2][]{%
\pgfkeys{/myfootnote, default, #1}%
\xdef\loc{\myfootnotePosition}%south east
\appear[\myappearCMD]{\begin{tikzpicture}[remember picture,overlay] 
		    \node (end) [yshift=-0.1cm,xshift=0.00cm,anchor=\loc] 
		    at (current page.\loc) {\subframetitlefont \tiny \myfootnoteRef};
		\end{tikzpicture}}%
}
	
\ifx\cleannotes\undefined
	%\newcommand{\appear}[1]{#1}
	\newcommand{\appear}[2][+-]{#2}%[2][+-]{#2}
	
	%% References for figures, for FYS.232 Structure of matter.
	\newcommand\HarrisCR[1][.-]{\tiny Figure from R. Harris, \textit{Modern Physics} (2$^{nd}$ edition).}
%	\newcommand\HarrisCR[1][.-]{\tiny \appear[#1]{Figure from R. Harris, \textit{Modern Physics} (2$^{nd}$ edition).}}
	\newcommand\HarrisCRshort[1][.-]{\tiny R. Harris, \textit{Modern Physics} (2$^{nd}$ edition).}
	\newcommand\TiplerCR[1][.-]{\tiny Figure from P. Tipler \& R. Llewellyn, \textit{Modern Physics} (5$^{th}$ edition).}
	\newcommand\TiplerCRshort[1][.-]{\tiny P. Tipler \& R. Llewellyn, \textit{Modern Physics} (5$^{th}$ edition).}
	\newcommand\CRTiplershort[1][.-]{\footnote{\TiplerCRshort[#1]}\vspace{-\baselineskip}}

	\NewDocumentCommand{\foocmd}{ O{default1} O{default2} m }{#1~#2~#3}
%	\NewDocumentCommand{\copyright}{ O{right} O{.-} m }{\footnote{\TiplerCRshort[#2]}\vspace{-\baselineskip}}
	\newcommand{\oldRightarrow}{\Rightarrow} % reduces problems with ctrl+f
	\newcommand{\oldrightarrow}{\rightarrow} % reduces problems with ctrl+f

	\def\plus{+}
	\def\equal{=}
	\def\minus{-}
	
	\newcommand{\Frac}[2]{\frac{#1}{#2}}
	\newcommand{\fracc}[2]{\frac{#1}{#2}}
	\newcommand{\fraca}[2]{\frac{#1}{#2}}
	%	\let\oldequiv\equiv
	%	\def\equiv{\appear{\oldequiv}}
	\def\Approx{\approx}
	\def\Equiv{\equiv}
	\def\[#1\]{%
		\vspace{-0.5mm}%
		  \begin{equation*} #1 \end{equation*}%
		  \vspace{-2.75mm}%
		  
		}
	\def\endofslide{\vspace{8cm} }
	\def\endofsection{\vspace{8cm} }	
\else
  \if\cleannotes1
	% \global pause to slide items
	\beamerdefaultoverlayspecification{<+->}
	% \appear cmd to pause arbitrary items
%	\newcommand{\appear}[1]{%
%	\uncover<+->{%
%	   #1}%
%	}%
	\newcommand{\appear}[2][+-]{%
	\uncover<#1>{%
	   #2}%
	}%
	% doesn't make sense since the _ and ^ markings won't be inside the same appear	
%	\renewcommand{\nsum}[1][1.4]{
%	  \mathop{%
%	    \raisebox
%	    {-#1\depthofsumsign+1.4\depthofsumsign}%-#1\depthofsumsign+1\depthofsumsign
%	    {\scalebox
%	      {#1}
%	      {$\displaystyle\sum$}%
%	    }
%	  }
%	}

	\newcommand\HarrisCR[1][.-]{\tiny \appear[#1]{Figure from R. Harris, \textit{Modern Physics} (2$^{nd}$ edition).}}
	\newcommand\HarrisCRshort[1][.-]{\tiny \appear[#1]{R. Harris, \textit{Modern Physics} (2$^{nd}$ edition).}}
	\newcommand\TiplerCR[1][.-]{\tiny \appear[#1]{Figure from P. Tipler \& R. Llewellyn, \textit{Modern Physics} (5$^{th}$ edition).}}
	\newcommand\TiplerCRshort[1][.-]{\tiny \appear[#1]{P. Tipler \& R. Llewellyn, \textit{Modern Physics} (5$^{th}$ edition).}}
	
	%\def\footnoterule{\only<3->\oldfootnoterule}
	
	%%%% Custom math operators that appear one by one
	\let\oldfrac\frac
	\newcommand{\Frac}[2]{\oldfrac{#1}{#2}}
%	\renewcommand{\frac}[2]{\appear{\oldfrac{#1}{#2}}}
%	\newcommand{\fracc}[2]{\appear{\oldfrac{\appear{#1}}{\appear{#2}}}}
%	\newcommand{\fraca}[2]{\appear{\oldfrac{\appear{#1}}{\appear{#2}}}}
%	\renewcommand{\frac}[2]{\appear{\oldfrac{#1}{#2}}}
	\renewcommand{\frac}[2]{{\appear{\oldfrac{#1}{#2}}}}
	\newcommand{\fracc}[2]{\appear{\oldfrac{\appear{#1}}{\appear{#2}}}}
	\newcommand{\fraca}[2]{\appear{\oldfrac{\appear{#1}}{\appear{#2}}}}
	
	\let\oldRightarrow\Rightarrow
	\let\oldrightarrow\rightarrow
	\renewcommand{\Rightarrow}{\appear{\oldRightarrow}} % reduces problems with ctrl+f
	\renewcommand{\rightarrow}{\appear{\oldrightarrow}} % reduces problems with ctrl+f

	%\gdef+{\appear{\textplus}}
%	\let\oldsqrt\sqrt
%	\def\sqrt{\appear{\oldsqrt}}

	\def\plus{\appear{+}}
	\def\equal{\appear{=}}
	\def\minus{\appear{-}}
	\let\oldequiv\equiv
	\def\equiv{\appear{\oldequiv}}
	\let\oldapprox\approx
	\def\approx{\appear{\oldapprox}}
	\def\Approx{\oldapprox}
	\def\Equiv{\oldequiv}
	\let\oldpropto\propto
	\renewcommand{\propto}{\appear{\oldpropto}}
	\newcommand{\Propto}{\oldpropto}
	%%
	%\def\[#1\]{%
	%  \begin{equation} \appear{#1} \end{equation}%
	%}
	\def\[#1\]{%
	  \vspace{-0.5mm}%
	  \begin{equation*} \appear{#1} \end{equation*}%
	  \vspace{-2.75mm}%
	  
	   }
	%\def$$#1$${%
	%  \begin{equation} \appear{#1} \end{equation}%
	%}
	%\newenvironment{displaymath}{begdef}{enddef}
	
	%\DeclareMathSymbol{+}{\mathbin}{operators}{43}
	%\begingroup
	%\catcode`+=\active
	%\gdef+{\mathbin{\text{\textplus}}}
	%\endgroup
	%\AtBeginDocument{\mathcode`+="8000 }
	
	\def\endofslide{\vspace{8cm}
		\appear{\begin{tikzpicture}[remember picture,overlay] 
	    \node (end) [yshift=-0.15cm,xshift=0.1cm,anchor=south east] 
	    at (current page.south east) {\oldrightarrow};
	\end{tikzpicture}%
	}%
	}
	\def\endofsection{\vspace{8cm}
		\appear{\begin{tikzpicture}[remember picture,overlay] 
		    \node (end) [yshift=-0.125cm,xshift=0.1cm,anchor=south east] 
		    at (current page.south east) {$\blacksquare$};
		\end{tikzpicture}%
		}
	}
   
  \fi
\fi

%\usepackage{amsmath,stackengine,mathtools}
%\stackMath

%\newcommand\CommentUnderset[3][red]{
%\underset{%
%	\appear{\textcolor{#1}{
%    \begin{array}{c}
%      #2
%    \end{array} }}
%}%
%{\cutline{#1}{#3}}
%}

\newcommand{\CommentUnderset}[3][red]{%
 \renewcommand\useanchorwidth{T}%tightens math to normal-looking eq.
 \stackunder[-1mm]{\cutline{#1}{#2}}
 {\scriptstyle\Alert[#1]{#3} }
}

%\usepackage{calc}
%\newlength{\depthofsumsign}
%\setlength{\depthofsumsign}{\depthof{$\sum$}}
%\newlength{\totalheightofsumsign}
%\newlength{\heightanddepthofargument}
%
%% https://tex.stackexchange.com/questions/22773/making-a-big-summation-sign
%\newcommand{\nsum}[1][1.4]{
%  \mathop{%
%    \raisebox
%    {-#1\depthofsumsign+1\depthofsumsign}
%    {\scalebox
%      {#1}
%      {$\displaystyle\sum$}%
%    }
%  }
%}
%
%\newcommand{\nint}[1][1.4]{
%  \mathop{%
%    \raisebox
%    {-#1\depthofsumsign+1\depthofsumsign}
%    {\scalebox
%      {#1}
%      {$\displaystyle\int$}%
%    }
%  }
%}
%
%\newcommand{\niint}[1][1.4]{
%  \mathop{%
%    \raisebox
%    {-#1\depthofsumsign+1\depthofsumsign}
%    {\scalebox
%      {#1}
%      {$\displaystyle\iint$}%
%    }
%  }
%}
%
%

%%%% GREEK SYMBOLS using unicode coding %%%%
\renewcommand\alpha{α} 
\renewcommand\beta{β} 

\renewcommand\gamma{γ} 
\renewcommand\Gamma{Γ} 

\renewcommand\delta{δ}
\renewcommand\Delta{Δ}

\renewcommand\epsilon{\symbol{"03B5}} %\symbol{"03B5} ϵ
\renewcommand\varepsilon{\symbol{"03B5}} %\symbol{"03B5} ϵ
\renewcommand\Epsilon{E}

\renewcommand\zeta{ζ}
\renewcommand\Zeta{Z}

\renewcommand\eta{η}
\renewcommand\Eta{H}

\renewcommand\theta{θ}
\renewcommand\Theta{Θ}

\renewcommand\iota{ι} 
\renewcommand\Iota{I}

\renewcommand\kappa{κ}
\renewcommand\Kappa{K}

\renewcommand\lambda{λ} 
\renewcommand\Lambda{Λ} 

\renewcommand\mu{μ} 
\renewcommand\Mu{M}

\renewcommand\nu{ν} 
\renewcommand\Nu{N}

\renewcommand\xi{ξ} 
\renewcommand\Xi{Ξ}

\renewcommand\xi{ξ} 
\renewcommand\Xi{Ξ}

\renewcommand\omicron{ξ} 
\renewcommand\Omicron{Ξ}

\renewcommand\pi{π} 
\renewcommand\Pi{Π}
 
%\renewcommand\rho{ξ} 
\renewcommand\rho{ρ} %ξ 
\renewcommand\Rho{Ξ}

\renewcommand\sigma{σ} 
\renewcommand\Sigma{Σ}

\renewcommand\tau{τ} 
\renewcommand\Tau{T}

\renewcommand\upsilon{υ} 
\renewcommand\Upsilon{ϒ}

\renewcommand\phi{ϕ} 
\renewcommand\Phi{Φ}

\renewcommand\chi{χ} 
\renewcommand\Chi{X}

\renewcommand\psi{ψ} 
\renewcommand\Psi{Ψ}

\renewcommand\omega{ω} 
\renewcommand\Omega{Ω}

\renewcommand\square{□}
\renewcommand\blacksquare{■}
\newcommand\blackdot{●}
\renewcommand\diamond{⋄}
\renewcommand\star{∗}

    
\renewcommand\hbar{ℏ} %ħ
 


